\chapter{Value, qualité et faible risque}\label{ch:s1:value-quality-and-low-risk}

À la fin de 2006, le facteur value américain avait fait passer un dollar investi en 1963 à 9.77~\$. En septembre 2020, il avait perdu 57.8~\% par rapport à ce sommet, soit presque deux
tiers de l'ensemble de ses gains en quarante-trois ans, et en juillet 2026 il était encore inférieur de 29.6~\% à ce niveau. Sur les mêmes années, le \omterm{def:s1:value-quality-and-low-risk:quality}{facteur de rentabilité} a continué de rapporter 2.8~\% par
an, avec un ratio de Sharpe de 0.51, et les actions à faible bêta ont continué de rapporter à peu près autant que les actions à bêta élevé. La value, la qualité et le faible risque forment le canon des
facteurs actions~: les plus étudiés, les plus échangés, les plus discutés. Ce chapitre les construit à partir de données point-in-time, montre que les détails de leur mise en œuvre
déterminent le résultat, et en examine le bilan. La construction s'appuie sur \texttt{firm.factorlib}.

\section{Value~: book-to-price et ses cousins}

\begin{definition}[Ratio valeur comptable sur cours, stratégie value]\label{def:s1:value-quality-and-low-risk:value}
Le \emph{ratio valeur comptable sur cours}\index{ratio valeur comptable sur cours} (book-to-price) d'une action est le rapport entre ses capitaux propres comptables par action et son cours (de manière équivalente, les capitaux propres comptables rapportés à la
capitalisation boursière). Une \emph{stratégie value}\index{stratégie value} achète les actions dont la mesure fondamentale (capitaux propres comptables, bénéfices, flux de trésorerie, chiffre d'affaires) rapportée au cours est élevée
et vend celles où ce ratio est faible.
\end{definition}

L'étude de Fama et French de 1992 sur la coupe transversale des rendements espérés constitue la référence standard pour la value, et leur facteur value, HML, est la moyenne des portefeuilles value de petites et de grandes capitalisations, moins la moyenne des portefeuilles growth de petites et de grandes capitalisations, reconstruit chaque mois de juin avec
les capitaux propres comptables de l'exercice précédent et la capitalisation boursière de décembre précédent. Sur la période de juillet 1963 à juillet 2026, il a rapporté 3.6~\% par an, avec un ratio de Sharpe
de 0.35 ($t = 2.77$)~; mais 5.7~\% par an jusqu'en 2006, $-5.3\%$ sur 2007--2020, et 9.4~\% depuis 2021 (la \cref{fig:s1:value-quality-and-low-risk:rolling}).

\begin{omfigure}
\begin{tikzpicture}
\begin{axis}[width=11cm,height=5.2cm,xlabel={\small année},ylabel={\small moyenne sur 10 ans (\% par an)},tick label style={font=\footnotesize},
  xmin=1972,xmax=2026,grid=major,grid style={gray!25},legend pos=south west,legend style={font=\scriptsize},table/col sep=comma,
  x tick label style={/pgf/number format/1000 sep={}}]
\addplot[omDef,very thick] table[x=year,y=hml_pct]{figdata/strategies-1/06-value-quality-and-low-risk/rolling.csv};
\addplot[omThm,thick] table[x=year,y=rmw_pct]{figdata/strategies-1/06-value-quality-and-low-risk/rolling.csv};
\addplot[black!50,dashed] coordinates {(1972,0) (2026,0)};
\legend{value (HML),rentabilité (RMW)}
\end{axis}
\end{tikzpicture}

\omcaption{Les facteurs value et rentabilité de Fama et French~: rendement moyen glissant sur 120 mois, annualisé, à chaque décembre. Calculé à partir de la Kenneth R. French
Data Library (Fama/French 5 factors, 2x3, mensuel, fichier CRSP 202607)~; la série brute n'est pas
redistribuée.}\label{fig:s1:value-quality-and-low-risk:rolling}
\end{omfigure}

\section{Qualité et rentabilité}

\begin{definition}[Facteur de rentabilité, facteur qualité]\label{def:s1:value-quality-and-low-risk:quality}
Un \emph{facteur de rentabilité}\index{facteur de rentabilité} est une position longue sur les actions d'entreprises à forte rentabilité (bénéfices bruts ou bénéfices d'exploitation rapportés aux actifs ou
aux capitaux propres comptables) et courte sur celles à faible rentabilité. Un \emph{facteur qualité}\index{facteur qualité} généralise ce principe à un composite de caractéristiques pour lesquelles les investisseurs
devraient payer davantage, comme la rentabilité, la croissance et la sécurité.
\end{definition}

Novy-Marx a montré que les bénéfices bruts rapportés aux actifs prédisent les rendements moyens à peu près aussi bien que le ratio valeur comptable sur capitalisation (book-to-market), que les entreprises rentables gagnent davantage malgré des
valorisations plus élevées, et que contrôler la rentabilité améliore nettement les \omterm{def:s1:value-quality-and-low-risk:value}{stratégies value}, surtout parmi les grandes capitalisations. Le modèle à cinq facteurs de Fama et French
(2015) ajoute la rentabilité d'exploitation (RMW, robuste moins faible) et l'investissement (CMA, conservateur moins agressif). Asness, Frazzini et Pedersen ont défini
la qualité comme la rentabilité, la croissance et la sécurité, et ont constaté que les actions de haute qualité se traitent à des cours plus élevés, mais pas de beaucoup, de sorte qu'un facteur quality-minus-junk
rapporte des rendements significatifs ajustés du risque aux États-Unis et dans 24 autres pays.

L'intérêt de ce couple tient aux corrélations. De 1963 à 2026, HML et RMW étaient corrélés à 0.09~; sur les années de pertes de la value, 2007--2020, $-0.14$. La value et
momentum étaient corrélés à $-0.19$, et à $-0.48$ sur 2007--2020. La value et l'investissement, en revanche, sont des cousins proches (0.68). La rentabilité a conservé sa prime
pendant la traversée de désert de la value~: 2.8~\% par an et un ratio de Sharpe de 0.51 sur 2007--2020.

Sur le marché synthétique, qui n'injecte aucune prime de rentabilité, un portefeuille de rentabilité (quatre fois le dernier bénéfice trimestriel rapporté aux capitaux propres comptables, tel que déposé) perd
3.8~\% par an, avec un ratio de Sharpe de $-1.03$~: c'est en réalité un portefeuille value vendeur déguisé (sa corrélation avec le portefeuille value est $-0.60$, et le coefficient d'information de son signal face à la
prime value injectée $-0.29$), puisqu'une valeur comptable élevée au dénominateur rend le ratio faible. Une caractéristique mesurée avec le même chiffre comptable que
un autre facteur hérite de la prime de ce facteur, avec le signe opposé.

\section{Faible risque et betting against beta}

\begin{definition}[Anomalie du faible risque, betting against beta]\label{def:s1:value-quality-and-low-risk:bab}
L'\emph{anomalie du faible risque}\index{anomalie du faible risque} est le constat que les actions à faible bêta et à faible volatilité ont rapporté des rendements ajustés du risque plus élevés que
les actions à bêta et à volatilité élevés~: la droite de marché des titres est plus plate que ce que prédit le modèle d'évaluation des actifs financiers (MEDAF). \emph{Betting against
beta}\index{betting against beta} (BAB) exploite cette anomalie au moyen d'un portefeuille long d'actions à faible bêta, avec un effet de levier qui porte leur bêta à un, et court d'actions à bêta élevé, avec un effet de levier qui ramène leur bêta à
un.
\end{definition}

L'explication de Frazzini et Pedersen est l'aversion au levier~: les investisseurs qui ne peuvent pas ou ne veulent pas emprunter achètent des actifs à bêta élevé pour s'exposer davantage au marché, en faisant
monter leur prix, de sorte qu'un bêta élevé va de pair avec un alpha faible. Ils l'ont observé sur les actions américaines, vingt autres marchés d'actions, les Treasuries, les obligations d'entreprise et les contrats à terme, et
ont constaté que les rendements du BAB sont faibles lorsque les contraintes de financement se resserrent. Les portefeuilles de French triés sur le bêta affichent la droite plate~: le décile au bêta le plus faible, avec un bêta réalisé
de 0.59, a rapporté 6.8~\% par an au-dessus du bon du Trésor (T-bill)~; le décile le plus élevé, avec un bêta de 1.60, a rapporté 8.6~\%. Le MEDAF, au rendement en excès du marché de 7.2~\%,
prédit 4.3~\% et 11.5~\%. Un portefeuille \omterm{def:s1:value-quality-and-low-risk:bab}{betting against beta} construit à partir des quintiles extrêmes, avec des bêtas glissants sur 60 mois, a rapporté 4.3~\% par an
depuis 1968, ratio de Sharpe de 0.26 ($t = 1.95$), avec un bêta de $-0.06$.

\begin{omfigure}
\begin{tikzpicture}
\begin{axis}[name=a,width=6.6cm,height=5.4cm,xlabel={\small bêta réalisé},ylabel={\small rendement en excès (\% par an)},tick label style={font=\footnotesize},
  xmin=0.4,xmax=1.7,ymin=0,ymax=16,grid=major,grid style={gray!25},title={\small déciles de bêta de French, 1963--2026},table/col sep=comma]
\addplot[omDef,only marks,mark=*] table[x=beta,y=ret_pct]{figdata/strategies-1/06-value-quality-and-low-risk/sml_french.csv};
\addplot[black!50,dashed,domain=0.4:1.7] {7.19*x};
\end{axis}
\begin{axis}[at={(a.east)},anchor=west,xshift=1.2cm,width=6.6cm,height=5.4cm,xlabel={\small bêta réalisé},tick label style={font=\footnotesize},
  xmin=0.4,xmax=1.7,ymin=0,ymax=16,grid=major,grid style={gray!25},title={\small marché synthétique, années 3--10},table/col sep=comma]
\addplot[omThm,only marks,mark=square*] table[x=beta,y=ret_pct]{figdata/strategies-1/06-value-quality-and-low-risk/sml_synth.csv};
\end{axis}
\end{tikzpicture}

\omcaption{Bêta réalisé et rendement moyen des déciles triés sur le bêta. À gauche~: les portefeuilles pondérés par la capitalisation de Kenneth French constitués sur le bêta, rendements en excès par rapport aux
bons du Trésor, la droite implicite du MEDAF au rendement en excès du marché de 7.2~\% (statistiques dérivées~; la série brute n'est pas redistribuée).
À droite~: les déciles de bêta glissant du marché synthétique, rendements totaux. Données~: \texttt{s1\_fetch\_value.py}, \texttt{s1\_factors.sml}.}\label{fig:s1:value-quality-and-low-risk:sml}
\end{omfigure}

Le marché synthétique a été construit à l'inverse~: ses rendements espérés valent le bêta multiplié par la prime de marché, plus les anomalies injectées, et rien ne récompense le bêta faible.
Sa droite de marché des titres est pentue (le décile au bêta glissant le plus faible a rapporté 2.4~\% par an à un bêta réalisé de 0.49, le plus élevé 13.4~\% à 1.39), et un
portefeuille \omterm{def:s1:value-quality-and-low-risk:bab}{betting against beta} perd 5.4~\% par an. Une stratégie dont toute la prémisse repose sur une droite de marché des titres plate n'a rien à gagner là où la droite n'est pas plate.

\section{Détails de mise en œuvre qui déterminent le résultat}

\begin{center}
\small
\begin{tabular}{@{}lcccc@{}}
\toprule
portefeuilles value synthétiques, années 3 à 10 & Sharpe & rendement & bêta & rotation\\
\midrule
portefeuille déposé, cours du jour, mensuel & $-0.07$ & $-0.7~\%$ & $-0.11$ & 3.2\\
portefeuille déposé, cours du jour, annuel & 0.21 & 1.9~\% & $-0.12$ & 2.5\\
timing Fama--French, annuel & 0.24 & 2.2~\% & $-0.12$ & 2.5\\
portefeuille look-ahead à partir de la fin de son exercice, mensuel & $-0.06$ & $-0.6~\%$ & $-0.11$ & 3.2\\
composite value et momentum, mensuel & 0.44 & 4.0~\% & 0.01 & 5.1\\
\bottomrule
\end{tabular}
\end{center}

La même prime de value implantée donne des ratios de Sharpe allant de $-0.07$ à 0.24 selon la date de référence du \omterm{def:s1:value-quality-and-low-risk:value}{ratio valeur comptable sur cours} (book-to-price). Le signal le plus précis (le cours du jour,
dont la corrélation de rang avec la prime implantée vaut 1.00) fait le moins bien, car le cours du jour intègre le rendement récent~: une action qui vient de chuter paraît
bon marché, et un portefeuille reconstruit chaque mois sur le cours du jour est vendeur du momentum implanté (sa corrélation avec le portefeuille momentum vaut $-0.56$). La convention de Fama--French
(un portefeuille âgé d'au moins six mois, divisé par le cours de sa propre date) porte moins de ce pari momentum et fait le mieux seul. Asness et Frazzini avancent le même
argument sur données réelles, mais par l'autre bout~: des cours actualisés prévoient mieux le vrai \omterm{def:s1:value-quality-and-low-risk:value}{ratio valeur comptable sur cours}, et les portefeuilles value construits sur eux ont dégagé des alphas de 3.05 à 3.78
points de pourcentage par an contre la méthode standard. Sur le marché synthétique, c'est la combinaison qui paie~: un composite de
value au cours du jour et momentum atteint un ratio de Sharpe de 0.44, meilleur que chacun des deux. Le portefeuille look-ahead, qui utilise chaque valeur comptable avant son dépôt, ne gagne
rien ici, car le portefeuille synthétique évolue lentement~; sur données réelles, un look-ahead de ce type est la façon la plus courante pour un backtest value de s'embellir.

\section{Combiner le canon et le timing des facteurs}

\begin{definition}[Timing des facteurs]\label{def:s1:value-quality-and-low-risk:timing}
Le \emph{timing des facteurs}\index{timing des facteurs} fait varier dans le temps les poids des facteurs d'un portefeuille selon des prévisions de leurs rendements, par exemple à partir de
le spread de valorisation entre le côté long et le côté short d'un facteur, sa performance récente ou sa volatilité.
\end{definition}

L'intérêt de combiner le canon tient à la diversification~: value, momentum et \omterm{def:s1:value-quality-and-low-risk:quality}{facteur de rentabilité} ont chacun connu des décennies perdantes, et leurs corrélations sont faibles ou négatives.
L'intérêt de faire du timing tient à la tentation du spread value~: quand les titres bon marché sont anormalement décotés par rapport aux titres chers, le rendement futur de la value
devrait être plus élevé. Asness, Frazzini et Pedersen en rapportent la version qualité~: un faible prix de la qualité annonçait un rendement futur élevé du QMJ. La
difficulté tient à ce que les signaux de timing sont lents, leurs échantillons courts, et que le bilan d'une décennie comme 2007--2020, où la value est restée bon marché et a continué de perdre, est la
raison pour laquelle la plupart des praticiens gardent le canon dans des proportions quasi fixes et le dimensionnent par la volatilité (chapitre~5) plutôt que d'en faire du timing.

\section{Fiches de stratégie}

\begin{strategyfile}{Value au book-to-price}\label{strat:s1:value-quality-and-low-risk:value}
\sfield{Qui paie, et pourquoi} Des investisseurs qui surréagissent aux mauvaises nouvelles et paient trop cher la croissance, ou une prime de risque pour la détresse~; la littérature n'est pas d'accord.
\sfield{Instruments et plateformes} Actions, en long sur les titres bon marché et en short sur les titres chers, ou inclinaisons long only.
\sfield{Signal} Capitaux propres comptables rapportés à la capitalisation boursière, sur données point-in-time.
\sfield{Dimensionnement et exécution} Tris neutres en taille (le 2x3 de HML) ou pondérations par caractéristique~; reconstruits chaque année ou chaque mois.
\sfield{Coûts} Faible rotation~: deux à trois fois le portefeuille par an ici.
\sfield{Comment elle s'éteint} De longues disettes (2007--2020)~; marchés tirés par la croissance~; comptabilité qui ne mesure plus la valeur (actifs incorporels).
\sfield{Horizon, capacité, infrastructure} Années~; capacité très grande~; fondamentaux point-in-time.
\sfield{Backtest en toute honnêteté} Valeurs comptables à partir de leurs dates de dépôt~; calage temporel du cours explicité~; la décennie perdue dans l'échantillon.
\sfield{Sources} Fama et French (1992)~; données de French~: HML 3.6~\% par an sur 1963--2026, ratio de Sharpe 0.35, un drawdown de 57.8~\% de décembre 2006 à septembre 2020~;
Asness et Frazzini (2013).
\end{strategyfile}

\begin{strategyfile}{Composite value}\label{strat:s1:value-quality-and-low-risk:composite}
\sfield{Qui paie, et pourquoi} Comme pour la value, l'erreur de mesure de chaque ratio pris isolément étant moyennée.
\sfield{Instruments et plateformes} Comme pour la value.
\sfield{Signal} La moyenne des z-scores de plusieurs ratios au cours (valeur comptable, bénéfice, flux de trésorerie, chiffre d'affaires), souvent accompagnée du momentum.
\sfield{Dimensionnement et exécution} Tris ou pondérations par caractéristique sur le composite.
\sfield{Coûts} Comme pour la value~; davantage avec le momentum dans le composite.
\sfield{Comment elle s'éteint} Comme pour la value.
\sfield{Horizon, capacité, infrastructure} Années~; plusieurs fondamentaux, point-in-time.
\sfield{Backtest en toute honnêteté} Les ingrédients du composite fixés à l'avance, et non choisis sur le backtest.
\sfield{Sources} La simulation de ce chapitre (value avec momentum~: ratio de Sharpe 0.44 contre 0.24 pour la value seule)~; Asness et Frazzini (2013).
\end{strategyfile}

\begin{strategyfile}{Rentabilité brute}\label{strat:s1:value-quality-and-low-risk:profitability}
\sfield{Qui paie, et pourquoi} Des investisseurs qui sous-évaluent la rentabilité durable, ou une prime que la littérature n'a pas tranchée.
\sfield{Instruments et plateformes} Actions, en long sur les sociétés rentables et en short sur les sociétés non rentables.
\sfield{Signal} Bénéfice brut (chiffre d'affaires moins coût des ventes) rapporté au total de l'actif.
\sfield{Dimensionnement et exécution} Tris, souvent combinés avec la value.
\sfield{Coûts} Faible rotation.
\sfield{Comment elle s'éteint} Encombrement (crowding)~; un dénominateur partagé avec la value, qui en fait un pari value déguisé.
\sfield{Horizon, capacité, infrastructure} Années~; capacité élevée~; comptes de résultat point-in-time.
\sfield{Backtest en toute honnêteté} Données comptables à partir des dates de dépôt~; sa corrélation avec les autres facteurs du portefeuille.
\sfield{Sources} Novy-Marx (2013)~; données de French~: RMW 3.1~\% par an sur 1963--2026, ratio de Sharpe 0.39.
\end{strategyfile}

\begin{strategyfile}{Quality minus junk}\label{strat:s1:value-quality-and-low-risk:qmj}
\sfield{Qui paie, et pourquoi} Des investisseurs qui paient trop peu cher les sociétés sûres, rentables et en croissance.
\sfield{Instruments et plateformes} Actions de nombreux pays.
\sfield{Signal} Un composite de scores de rentabilité, de croissance et de sécurité.
\sfield{Dimensionnement et exécution} En long sur la qualité, en short sur le junk~; neutre en taille et en secteur.
\sfield{Coûts} Faible rotation.
\sfield{Comment elle s'éteint} Quand la qualité devient chère~; le prix de la qualité prédit son rendement.
\sfield{Horizon, capacité, infrastructure} Années~; nombreuses données comptables en entrée.
\sfield{Backtest en toute honnêteté} Chaque donnée en entrée point-in-time~; recette du composite fixée à l'avance.
\sfield{Sources} Asness, Frazzini et Pedersen (2019)~: rendements ajustés du risque significatifs aux États-Unis et dans 24 autres pays.
\end{strategyfile}

\begin{strategyfile}{Betting against beta}\label{strat:s1:value-quality-and-low-risk:bab}
\sfield{Qui paie, et pourquoi} Des investisseurs dont le levier est contraint, qui achètent des actifs à bêta élevé plutôt que d'emprunter.
\sfield{Instruments et plateformes} Actions, obligations et contrats à terme, avec effet de levier du côté long.
\sfield{Signal} Bêta estimé.
\sfield{Dimensionnement et exécution} En long sur les bêtas faibles, ajustés à un~; en short sur les bêtas élevés, ajustés à un.
\sfield{Coûts} Financement du côté long à effet de levier~; vente à découvert des titres à bêta élevé.
\sfield{Comment elle s'éteint} Resserrement du financement (les rendements du BAB sont alors faibles)~; droite de marché des titres très pentue.
\sfield{Horizon, capacité, infrastructure} Mois~; estimations de bêta, financement.
\sfield{Backtest en toute honnêteté} Bêtas connus avant chaque mois~; coûts de financement~; effet de levier du côté long.
\sfield{Sources} Frazzini et Pedersen (2014)~; portefeuilles de bêta de French~: 4.3~\% par an depuis 1968, ratio de Sharpe 0.26, avant coûts de financement et coûts de transaction.
\end{strategyfile}

\begin{strategyfile}{Faible volatilité long only}\label{strat:s1:value-quality-and-low-risk:lowvol}
\sfield{Qui paie, et pourquoi} Comme pour le BAB, pris en portefeuille long only avec moins de risque de marché.
\sfield{Instruments et plateformes} Actions, long only.
\sfield{Signal} Bêta faible ou volatilité faible.
\sfield{Dimensionnement et exécution} Le quintile le moins risqué ou un portefeuille de variance minimale (livre~7, chapitre~26).
\sfield{Coûts} Faible rotation.
\sfield{Comment elle s'éteint} Encombrement sur les mêmes valeurs défensives~; hausse des taux qui frappe les actions assimilables à des obligations.
\sfield{Horizon, capacité, infrastructure} Années~; un modèle de risque.
\sfield{Backtest en toute honnêteté} On le juge sur le rendement ajusté du risque et face au marché, avec son propre bêta.
\sfield{Sources} Portefeuilles de bêta de French~: le décile le plus bas a rapporté 6.8~\% par an au-dessus des bons du Trésor, pour un bêta de 0.59, sur 1963--2026.
\end{strategyfile}

\section{Tutoriel~: la décennie perdue de la value}\label{tut:s1:value-quality-and-low-risk:tut}

\begin{tutorial}
\textbf{Objectif.} Construire les portefeuilles value, de rentabilité et \omterm{def:s1:value-quality-and-low-risk:bab}{betting against beta} à partir de fondamentaux point-in-time, dater le \omterm{def:s1:value-quality-and-low-risk:value}{ratio valeur comptable sur cours} de quatre façons, et lire l'historique
de French. \textbf{État final~:} le tableau d'implémentation et la \cref{fig:s1:value-quality-and-low-risk:sml}.
\begin{tutsteps}
\item \textbf{Book-to-price calé dans le temps}~: avec le cours du jour ou avec le cours à sa propre date.
\omcode{code/firm/factorlib/firm_factorlib.py}{86}{96}{Book-to-price à partir des valeurs comptables déposées.}\label{lst:s1:value-quality-and-low-risk:bp}
\item \textbf{\omterm{def:s1:value-quality-and-low-risk:bab}{Betting against beta}}~: chaque côté ramené à un bêta de un.
\omcode{code/firm/factorlib/firm_factorlib.py}{70}{83}{Le portefeuille betting against beta.}\label{lst:s1:value-quality-and-low-risk:bab}
\item \textbf{Lancer} \texttt{s1\_factors.run(name)} pour les sept portefeuilles, \texttt{sml()}, \texttt{correlations()}, puis \texttt{s1\_fetch\_value.py} une fois, et
\texttt{fig\_factors.py}.
\end{tutsteps}
\textbf{Pour aller plus loin.} Pondérer les portefeuilles par capitalisation boursière (l'argument \texttt{cap} de \texttt{sort\_book})~; ajouter le spread value comme signal de timing et le tester sur l'échantillon de French
avant et après 2007~; implanter une prime de rentabilité et vérifier que le portefeuille de rentabilité la retrouve.
\end{tutorial}

\section{Construction~: la bibliothèque de facteurs}\label{bld:s1:value-quality-and-low-risk:factorlib}

\begin{build}
\bfield{Rôle} Le canon des facteurs actions, construit à partir de caractéristiques point-in-time, sous les formes qu'utilisent les praticiens et les articles.
\bfield{Interface} \texttt{sort\_book(signal, universe, q, cap)}, \texttt{char\_book(signal, universe)}, \texttt{composite(*signals)},
\texttt{bab\_book(beta, universe, q)}, \texttt{book\_to\_price(anchor, cum, listed, current)}.
\bfield{Règles} Caractéristiques établies à partir de leurs dates de dépôt~; calage temporel des prix explicite~; chaque facteur présenté avec ses corrélations aux autres.
\bfield{Tests d'acceptation} \texttt{code/firm/factorlib/tests/}~: pondérations des tris et des caractéristiques vérifiées à la main~; composites avec valeurs manquantes~; bêta nul du BAB~;
book-to-price avec le cours du jour et le cours calé.
\bfield{Pour aller plus loin} Tris 2x3 neutres en taille, avec points de rupture du NYSE~; rentabilité brute à partir des comptes de résultat (livre~7, chapitre~11, \texttt{firm.fundpit})~;
timing du spread value.
\end{build}

\begin{omsources}
\item E.~F.~Fama and K.~R.~French, ``The cross-section of expected stock returns'', \emph{Journal of Finance} 47(2), 1992.
\item R.~Novy-Marx, ``The other side of value: the gross profitability premium'', \emph{Journal of Financial Economics} 108(1), 2013.
\item C.~S.~Asness, A.~Frazzini and L.~H.~Pedersen, ``Quality minus junk'', \emph{Review of Accounting Studies} 24(1), 2019.
\item A.~Frazzini and L.~H.~Pedersen, ``Betting against beta'', \emph{Journal of Financial Economics} 111(1), 2014.
\item E.~F.~Fama and K.~R.~French, ``A five-factor asset pricing model'', \emph{Journal of Financial Economics} 116(1), 2015.
\item C.~Asness and A.~Frazzini, ``The devil in HML's details'', \emph{Journal of Portfolio Management} 39(4), 2013.
\item Kenneth R. French Data Library~: Fama/French 5 factors (2x3), facteur momentum, portefeuilles constitués sur le bêta.
\end{omsources}

\section{Exercices}

\begin{exercise}[$\star$]\label{exo:s1:value-quality-and-low-risk:1}
La richesse cumulée de HML valait 9.77~\$ par dollar à son sommet de décembre 2006, puis a chuté de 57.8~\%. Quelle part de son gain cumulé a-t-elle perdue~?
\end{exercise}

\begin{exercise}[$\star$]\label{exo:s1:value-quality-and-low-risk:2}
Le côté à faible bêta d'un portefeuille BAB a un bêta de 0.6 et son côté à fort bêta un bêta de 1.5. Combien de dollars le portefeuille détient-il en position longue et en position courte par dollar de bêta de chaque côté, et
quelle est sa position nette en dollars~?
\end{exercise}

\begin{exercise}[$\star$]\label{exo:s1:value-quality-and-low-risk:3}
Le marché synthétique verse 3~\% par an par écart type du book-to-price. Que rapporte un portefeuille par déciles si ses déciles longs et courts affichent en moyenne $\pm 1.75$
écarts types, avec un demi-dollar de chaque côté~?
\end{exercise}

\begin{exercise}[$\star\star$]\label{exo:s1:value-quality-and-low-risk:4}
Expliquer pourquoi le book-to-price calculé avec le cours du jour est négativement corrélé au momentum, et pourquoi cela fait de la \omterm{def:s1:value-quality-and-low-risk:value}{stratégie value} et du momentum de bons partenaires.
\end{exercise}

\begin{exercise}[$\star\star$]\label{exo:s1:value-quality-and-low-risk:5}
Pourquoi le portefeuille synthétique de rentabilité perd-il lorsqu'aucune prime de rentabilité n'est implantée~?
\end{exercise}

\begin{exercise}[$\star\star$]\label{exo:s1:value-quality-and-low-risk:6}
Pour un rendement excédentaire du marché de 7.2~\%, que prédit le modèle d'évaluation des actifs financiers (MEDAF) pour des bêtas de 0.594 et 1.604~? Comparer avec les déciles de French.
\end{exercise}

\begin{exercise}[$\star\star\star$]\label{exo:s1:value-quality-and-low-risk:7}
\emph{Programmation.} Lancer \texttt{run('composite')} et \texttt{correlations()}. Expliquer le ratio de Sharpe du composite à partir de la corrélation des deux portefeuilles.
\end{exercise}

\begin{exercise}[$\star\star\star$]\label{exo:s1:value-quality-and-low-risk:8}
\emph{Trouver la faille.} «~La \omterm{def:s1:value-quality-and-low-risk:value}{stratégie value} ne fonctionne plus depuis 2007, elle a donc été éliminée par l'arbitrage~; nous l'abandonnons.~»
\end{exercise}

\section{Problème~: La décennie perdue de la stratégie value}

\begin{problem}[{Problème du week-end --- le canon, construit et analysé}]\label{pb:s1:value-quality-and-low-risk:1}
Les données factorielles de French et les portefeuilles factoriels synthétiques du chapitre.

\medskip\noindent\textbf{Partie I --- Value.}
\begin{enumerate}
\item Définir le book-to-price et une \omterm{def:s1:value-quality-and-low-risk:value}{stratégie value}, et décrire la construction de HML.
\item Donner le rendement et le ratio de Sharpe de HML sur l'ensemble de l'échantillon et sur les trois périodes.
\item Décrire le drawdown de 2007--2020.
\item Pourquoi le calage temporel du prix fait-il partie de la définition~?
\end{enumerate}

\medskip\noindent\textbf{Partie II --- Qualité.}
\begin{enumerate}[resume]
\item Définir les \omterm{def:s1:value-quality-and-low-risk:quality}{facteurs de rentabilité} et de qualité et donner leurs résultats publiés.
\item Donner les corrélations de HML avec RMW, le momentum et CMA, sur l'ensemble de l'échantillon et en 2007--2020.
\item Quel rendement le \omterm{def:s1:value-quality-and-low-risk:quality}{facteur de rentabilité} a-t-il rapporté pendant la période de sécheresse de la value~?
\item Pourquoi le portefeuille synthétique de rentabilité perd-il~?
\end{enumerate}

\medskip\noindent\textbf{Partie III --- Faible risque.}
\begin{enumerate}[resume]
\item Définir l'\omterm{def:s1:value-quality-and-low-risk:bab}{anomalie du faible risque} et le \omterm{def:s1:value-quality-and-low-risk:bab}{betting against beta} (BAB).
\item Que montrent les données de French triées par bêta, face au MEDAF~?
\item Donner le rendement, le ratio de Sharpe et le bêta du portefeuille BAB.
\item Pourquoi le BAB perd-il sur le marché synthétique~?
\end{enumerate}

\medskip\noindent\textbf{Partie IV --- Le verdict.}
\begin{enumerate}[resume]
\item Énoncer le \emph{résultat nommé}~: le drawdown du facteur value depuis son sommet de 2007 dans les données de French, et la corrélation des rendements du facteur value et du \omterm{def:s1:value-quality-and-low-risk:quality}{facteur de rentabilité}.
\item Classer les quatre calages temporels du book-to-price synthétique et expliquer ce classement.
\item Que donne la combinaison de la value et du momentum ici, et pourquoi~?
\item Que trouvent Asness et Frazzini au sujet des prix correctement datés~?
\item Définir le \omterm{def:s1:value-quality-and-low-risk:timing}{timing des facteurs} et donner un argument pour et un argument contre.
\item Que faut-il vérifier avant de faire confiance à un backtest de la value~?
\item Comment répartir le capital entre les facteurs du canon~?
\item En une phrase, qu'ont en commun la value, la qualité et le faible risque~?
\end{enumerate}
\end{problem}

\section{Questions d'entretien}

\begin{interviewq}[$\star$ \iqroles{chercheur}]\label{iq:s1:value-quality-and-low-risk:1}
Comment le facteur value de Fama--French est-il construit~?
\end{interviewq}

\begin{interviewq}[$\star\star$ \iqroles{chercheur}]\label{iq:s1:value-quality-and-low-risk:2}
Pourquoi les actions à faible bêta pourraient-elles rapporter des rendements ajustés du risque supérieurs à ceux des actions à bêta élevé~?
\end{interviewq}

\begin{interviewq}[$\star\star$ \iqroles{chercheur, risque}]\label{iq:s1:value-quality-and-low-risk:3}
La \omterm{def:s1:value-quality-and-low-risk:value}{stratégie value} a perdu plus de la moitié de sa valeur entre 2007 et 2020. Que concluriez-vous, et que ne concluriez-vous pas~?
\end{interviewq}

\begin{interviewq}[$\star\star$ \iqroles{chercheur}]\label{iq:s1:value-quality-and-low-risk:4}
Pourquoi combiner la \omterm{def:s1:value-quality-and-low-risk:value}{stratégie value} avec le \omterm{def:s1:value-quality-and-low-risk:quality}{facteur de rentabilité}, ou avec le momentum~?
\end{interviewq}

\begin{interviewq}[$\star\star$ \iqroles{développeur, chercheur}]\label{iq:s1:value-quality-and-low-risk:5}
Quels problèmes liés aux données point-in-time se posent lors de la construction d'un signal de book-to-price~?
\end{interviewq}

\begin{interviewq}[$\star\star\star$ \iqroles{chercheur}]\label{iq:s1:value-quality-and-low-risk:6}
Deux facteurs ont chacun un ratio de Sharpe de 0.3 et une corrélation de $-0.5$. Quel est le ratio de Sharpe de leur combinaison à pondérations égales, et celui de la meilleure
combinaison~?
\end{interviewq}
